\documentclass[11pt]{article}
\usepackage{icmlko}
\begin{document}

\title{자기가 풀 수 있는 문제만 채점받고 있었다 --- 기준선을 고치자 우리 우위가 세 배가 되고, 강조가 왜 아무것도 못 고쳤는지가 설명됐다}
\author{월드모델 실험실}
\date{노트 529 · 2026년 8월 4일}
\maketitle

\begin{abstract}
노트 527 의 사전 등록 --- ``축 하나가 지배하는 도메인에서 공유 모델이
희석한다, 그러니 \textbf{도메인 전용 축 사본}(\texttt{DOMAX}, 가림의 거울)을
주면 오른다'' --- 을 쟀다. \textbf{반증이다.}

\begin{center}
\begin{tabular}{lrr}
\hline
팔 & 그 도메인 차 & 부호 \\
\hline
모바일 \texttt{entry\_friction} & $+0.0002$ & $4/8$ \\
시장팝업 \texttt{goods\_scale} & $-0.0030$ & $2/8$ \\
팝업 \texttt{target\_breadth} & $+0.0002$ & $2/8$ \\
\textbf{아이돌 \texttt{venue\_prominence}} & $\mathbf{+0.0525}$ & $\mathbf{8/8}$ \\
\hline
\end{tabular}
\end{center}

넷 중 \emph{하나}만 올랐고 판은 $+0.0006$($5/8$, $[-0.0007,+0.0020]$)로
$0$ 을 문다. \textbf{그런데 왜 모바일이 꿈쩍도 안 했는지를 파다가 기준선
자체가 틀렸다는 것을 찾았다.}
\end{abstract}

\begin{figure}[t]
\centering
\includegraphics[width=\textwidth]{figs/f1.pdf}
\caption{\textbf{가} --- 단일 축의 동점률. \textbf{나} --- 채점 쌍을 맞추면.
\textbf{다} --- 두 실험이 맞물린다.}
\end{figure}

\section{모바일 \texttt{entry\_friction} 은 앱 가격이다}

모바일에서 이 축 하나가 $83.2\%$ 로 우리($72.5\%$)를 $10.7$ \%p 앞선다고 적었다.
그런데 \texttt{provenance} 를 보면 모바일 \texttt{entry\_friction} 은
\textbf{앱 가격}이고 무료가 $0$ 이다 --- \emph{값이 거의 이산적}이다.

내 쌍 정확도 함수는 \textbf{예측값이 같은 쌍을 제외}한다. 모델 예측은
연속이라 동점이 거의 없지만 가격은 동점이 대량이다.

\begin{quote}
\textbf{모델은 전체 쌍에서 채점받고 단일 축은 자기가 구별하는 쌍에서만
채점받고 있었다.} 같은 판이 아니다.
\end{quote}

\begin{center}
\begin{tabular}{lrrr}
\hline
도메인 & 축(동점 뺌) & \textbf{동점률} & 우리(같은 쌍) \\
\hline
시장팝업 & $71.1\%$ & $\mathbf{89.4\%}$ & $71.1\%$ \\
도서 & $36.9\%$ & $\mathbf{92.9\%}$ & $58.8\%$ \\
게임 & $69.9\%$ & $81.6\%$ & $76.3\%$ \\
모바일 & $83.2\%$ & $68.9\%$ & $\mathbf{84.1\%}$ \\
아이돌 & $78.2\%$ & $49.4\%$ & $75.0\%$ \\
\hline
가중 평균 & $64.4\%$ & $\mathbf{37.2\%}$ & $\mathbf{70.6\%}$ \\
\hline
\end{tabular}
\end{center}

\section{고치면 우위가 세 배가 된다}

\begin{center}
\begin{tabular}{lrrr}
\hline
견주는 법 & 축 & 우리 & 차 \\
\hline
노트 527(축이 고른 쌍) & $64.4\%$ & $67.2\%$ & $+2.9$ \%p \\
\textbf{같은 쌍에서만} & $64.4\%$ & $\mathbf{70.6\%}$ & $\mathbf{+6.3}$ \textbf{\%p} \\
\textbf{동점을 $\frac{1}{2}$ 로}(정석) & $\mathbf{57.9\%}$ & $67.2\%$ & $\mathbf{+9.3}$ \textbf{\%p} \\
\hline
\end{tabular}
\end{center}

두 셈법 다 정당하고 답하는 물음이 다르다 --- ``축이 구별할 수 있을 때 누가
낫나''(같은 쌍)와 ``전체 인구에서 누가 낫나''(동점 $\frac{1}{2}$, AUC $\cdot$
$\tau_b$ 의 관례). \textbf{둘 다 우리 쪽이 노트 527 보다 훨씬 낫다.}

\begin{quote}
\textbf{모바일에서 진다는 것은 착시였다} --- 같은 쌍에서 $84.1\%$ 대
$83.2\%$ 로 \emph{우리가 이긴다}. 시장팝업도 $71.1$ 대 $71.1$ 로 동률이다.
\textbf{같은 쌍에서 진짜로 지는 곳은 아이돌($75.0$ 대 $78.2$)과
팝업($61.6$ 대 $67.1$) 둘뿐}이고, 둘 다 유보가 $51 \cdot 65$ 행으로 제일 작다.
\end{quote}

\section{두 실험이 서로를 확인한다}

이것이 \texttt{DOMAX} 결과를 그대로 설명한다.

\begin{quote}
\textbf{강조가 크게 올린 유일한 도메인(아이돌 $+0.0525$, $8/8$)이 같은 쌍에서도
크게 지는 유일한 도메인이다.} 모바일이 꿈쩍 안 한 이유는 희석되고 있어서가
아니라 \textbf{애초에 고칠 것이 없어서}다.
\end{quote}

\emph{한 실험의 실패가 다른 실험의 오류를 잡았다.} 사전 등록이 반증됐을 때
``왜''를 판 것이 값을 냈다(노트 133).

\section*{조항}

\begin{quote}
\textbf{예측기를 견줄 때는 채점하는 쌍 집합을 맞춘다.} 이산적인 예측기는
자기가 구별 못 하는 쌍을 채점에서 빼면 유리해진다. 기본은 \textbf{동점을
$\frac{1}{2}$ 로} 세는 것이고, ``구별할 수 있을 때의 실력''을 보고 싶으면
\textbf{양쪽 다 같은 쌍}에서 잰다.
\end{quote}

이 조항은 우리 자신에게도 적용된다 --- 우리 모델은 예측이 연속이라 동점이
거의 없고, 그래서 이 함정에서 \emph{구조적으로 유리한} 쪽이었다. 유리한
쪽에 서 있을 때 규칙을 의심하기가 더 어렵다.

\section*{\texttt{DOMAX} 의 처분}

채택하지 않는다 --- 판 $+0.0006$ 이 $0$ 을 물고, 노트 507 조항($+0.002$ 아래는
씨앗 $40$ 전 채택 금지)에도 걸린다. 다만 \textbf{코드는 남긴다}
(\texttt{DOMAX} 기본 빈 사전 $=$ 챔피언 불변).

\begin{quote}
\textbf{아이돌 하나는 진짜다} --- 도메인 $+0.0525$ 에 $8/8$. 그런데 아이돌은
판 무게가 $1.5\%$ 라 \textbf{판이 그것을 거의 못 본다.} 노트 528 이
``관문의 자 다섯이 접전 개선을 못 본다''고 적은 것과 같은 모양이다 ---
\textbf{우리 관문이 우리가 원하는 개선을 못 보는 자리가 늘고 있다.}
\end{quote}

\section*{남은 것}

\begin{itemize}
\item \textbf{노트 527 의 수치를 원장에서 고쳐 적는다}(우위 $+3.0 \to +6.3$ /
$+9.3$ \%p, ``지는 넷''$\to$``지는 둘'').
\item \textbf{작은 도메인 전용 관문} --- 아이돌 $\cdot$ 팝업처럼 판 무게가
$2\%$ 아래인 곳의 개선을 어떻게 채택할지. 지금 구조로는 영영 못 채택한다.
\item \textbf{동점률이 높은 축은 축 자체가 거친 것} --- 시장팝업
\texttt{goods\_scale} $89.4\%$ $\cdot$ 도서 \texttt{grp} $92.9\%$.
태거를 고칠 자리를 가리키는 새 지표일 수 있다.
\end{itemize}

\end{document}
